\section{Lattice sequences and prefix-based CBC construction}
\label{sec:lattice}

Fix a maximum sample size $N\le 2^M$ and an integer generating vector
$\vh=(h_1,\ldots,h_d)$. The base-two radical inverse is
\[
 \phi_2(i)=\sum_{k=0}^{M-1}i_k2^{-k-1},
 \qquad i=\sum_{k=0}^{M-1}i_k2^k,\quad 0\le i<2^M.
\]
Our lattice sequence has points
$\vx_i=\fracpart{\phi_2(i)\vh+\vDelta}$, with one common uniform random
shift $\vDelta$. At a dyadic sample size $n=2^m$, its prefix is the usual
rank-one lattice point set with denominator $2^m$, provided the components of
$\vh$ are odd. At other sample sizes the prefix remains well defined, although
it need not be a full lattice point set. We search for a single generating
vector using all these prefixes.

\subsection{Counting scalar differences}

The shift-averaged formula~\eqref{eq:shift-esd} becomes
\begin{equation}\label{eq:lattice-esd}
 e_n^2(\vh)
 =-1+\frac{\kt_d(\vzero)}{n}
  +\frac{1}{n^2}\sum_{r=1}^{2^M-1}
       \rho_n(r)\kt_d\bigl(\fracpart{\phi_2(r)\vh}\bigr),
\end{equation}
where $\rho_n(r)$ counts the ordered pairs of distinct indices in the prefix
whose scalar difference, modulo one, is $\phi_2(r)$. The constant $-1$ uses
the normalization of~\eqref{eq:kernel}.
The count is constant on each dyadic block of radical-inverse indices.
For $2^{k-1}\le r<2^k$, $1\le k\le M$, it is
\begin{equation}\label{eq:rho}
 \rho_n(r)=r_{n,k}:=
 n-(n\bmod 2^k)
 +2\max\{(n\bmod 2^k)-2^{k-1},0\}.
\end{equation}
This formula and the aggregated criterion below follow the calculations in
Pride's lattice-search working notes~\cite{pride-lattice-notes}.

To see the count, add prefix indices $i=0,1,\ldots,n-1$ in order.
The differences from $\phi_2(i)$ to its predecessors contain all the odd
multiples of $2^{-k}$ exactly once if bit $k-1$ of $i$ equals one, and contain
none otherwise. The reverse ordered differences give the same collection.
Thus $r_{n,k}$ is twice the number of indices $i<n$ having that bit equal to
one. Counting complete cycles of length $2^k$ and then the remaining indices
gives~\eqref{eq:rho}. Section~\ref{sec:newo} proves the corresponding difference
identity for a larger class of scalar sequences.

Define
\begin{equation}\label{eq:lattice-blocks}
 A_k(\vh)=\sum_{r=2^{k-1}}^{2^k-1}
       \kt_d\bigl(\fracpart{\phi_2(r)\vh}\bigr).
\end{equation}
For arbitrary nonnegative sample-size weights $w_n$, equations
\eqref{eq:lattice-esd}--\eqref{eq:lattice-blocks} give
\begin{align}\label{eq:lattice-weighted}
 \sum_{n=1}^N w_ne_n^2(\vh)
 &=\kt_d(\vzero)\sum_{n=1}^N\frac{w_n}{n}
   -\sum_{n=1}^Nw_n
   +\sum_{k=1}^Mq_kA_k(\vh),\\
 q_k&=\sum_{n=1}^N\frac{w_n}{n^2}r_{n,k}.\notag
\end{align}
In particular, the criterion $\W_N$ in~\eqref{eq:wesd}, with $w_n=n$, is
\begin{equation}\label{eq:lattice-wesd}
 \W_N(\vh)=N\kt_d(\vzero)-\frac{N(N+1)}2
           +\sum_{k=1}^Mq_kA_k(\vh),
 \qquad q_k=\sum_{n=1}^N\frac{r_{n,k}}n.
\end{equation}
The $q_k$ depend on the sample-size range and weights, not on the generating
vector, so they can be computed once for an entire construction.

\subsection{A component-by-component search}

Write the product kernel~\eqref{eq:kernel} as
$\kt_s(\vx)=\prod_{j=1}^s\kappa_j(x_j)$, with the coordinate weights
$\gamma_j$ included in the one-dimensional factors $\kappa_j$.
At stage $s$, retain the previously selected $h_1,\ldots,h_{s-1}$ and cache
\[
 P_{s-1}(r)=\prod_{j=1}^{s-1}
       \kappa_j\bigl(\fracpart{\phi_2(r)h_j}\bigr),
 \qquad 1\le r<2^M.
\]
For a candidate $h$, evaluate
\begin{equation}\label{eq:lattice-cbc-score}
 Q_s(h)=\sum_{k=1}^Mq_k
       \sum_{r=2^{k-1}}^{2^k-1}
       P_{s-1}(r)\kappa_s\bigl(\fracpart{\phi_2(r)h}\bigr).
\end{equation}
The other terms of~\eqref{eq:lattice-weighted} do not depend on $h$.
Consequently minimizing $Q_s$ is exactly minimizing the weighted prefix
criterion for the current dimension. Choose a minimizer, break ties by a
fixed rule, and multiply the cache by the selected new factor before
proceeding to coordinate $s+1$.

A transparent finite search uses the odd residues
$\{1,3,\ldots,2^M-1\}$ as candidates. The first component may be fixed to one
in the conventional construction; all subsequent comparisons then use that
restriction. Each candidate score requires $O(2^M)$ operations once the
coefficients and cache are available. This describes a direct CBC method,
not an FFT-accelerated complexity claim. It is greedy across coordinates:
no assertion of a globally minimizing generating vector follows merely from
the componentwise minimizations.

\par\smallskip
\noindent\codex{[Author verification: Confirm the intended candidate restrictions and optimized implementation
with Anders. The formulas above specify a reproducible direct search;
computational speedups and their complexity still need to be aligned with the
implementation used for the final numerical comparisons.]}
\par\smallskip
